\section{Bài toán và cách biểu diễn dữ liệu}
Trích xuất thông tin chính (Key Information Extraction, KIE) là xác định vai trò của các đoạn chữ và đưa chúng về các trường dữ liệu. Ví dụ, từ một hóa đơn, hệ thống cần lấy đúng tên cửa hàng, địa chỉ, ngày mua và tổng tiền. Hai chuỗi số giống nhau có thể mang ý nghĩa khác nhau tùy vị trí và nội dung xung quanh.

Nhận dạng chữ (OCR) cung cấp văn bản và vùng bao chữ, còn KIE xác định ý nghĩa của chúng. Với các cấu hình cần OCR, văn bản và vùng bao được lấy từ chú thích gốc của bộ dữ liệu. Kết quả vì vậy phản ánh khả năng trích xuất thông tin khi đã có đầu vào chữ và vị trí, chưa đánh giá sai số của một hệ OCR thực tế.

Trong các mô hình gán nhãn, BIO quy định: B đánh dấu phần bắt đầu, I đánh dấu phần tiếp theo và O chỉ phần không thuộc trường cần lấy. Chẳng hạn, hai từ trong ``ABC STORE'' có thể nhận nhãn \texttt{B-company} và \texttt{I-company}. Chương trình ghép các phần được gán nhãn để tạo giá trị của trường \texttt{company}.

\section{Transformer và học chuyển giao}
Transformer tạo biểu diễn của một từ dựa trên các từ khác trong chuỗi thông qua cơ chế attention \parencite{vaswani2017}. Nhờ đó, một chuỗi số có thể được diễn giải dựa trên các từ như ``TOTAL'' hoặc ``TAX''. Tuy nhiên, thứ tự trong chuỗi chưa mô tả đầy đủ vị trí hai chiều trên hóa đơn.

Tiền huấn luyện giúp mô hình học biểu diễn từ một lượng dữ liệu lớn; fine-tuning tiếp tục điều chỉnh trọng số bằng dữ liệu có nhãn của bài toán cụ thể \parencite{devlin2019}. Nhóm sử dụng trọng số RoBERTa và BROS đã tiền huấn luyện rồi fine-tune cho nhãn BIO. Với PICK, nhánh ResNet-50 sử dụng trọng số tiền huấn luyện, các thành phần còn lại được khởi tạo mới.
